# 13. Vragen

Beantwoord de volgende vragen in je verslag.

### Vraag 1

Wat is het verschil tussen een **gebruiker** en een **groep** in Linux?

### Vraag 2

Waarom gebruik je groepen bij het delen van bestanden?

### Vraag 3

Wat doet het commando:

```bash
chmod
```

### Vraag 4

Wat doet:

```bash
chown
```

### Vraag 5

Wat doet:

```bash
chgrp
```

### Vraag 6

Wat betekent deze rechtenreeks?

```text
drwxrws---
```

### Vraag 7

Wat is het doel van het setgid-bit op een gedeelde directory?

### Vraag 8

Waarom is het niet voldoende om alleen de gebruiker als eigenaar van een bestand in te stellen wanneer meerdere medewerkers met hetzelfde bestand moeten werken?

### Vraag 9

Wat gebeurt er als een gebruiker geen rechten heeft op een directory, maar wel rechten heeft op een bestand dat zich in die directory bevindt?

### Vraag 10

Waarom is het vanuit beveiliging belangrijk dat gebruikers alleen toegang krijgen tot de mappen die zij nodig hebben?

---

# 14. Eindopdracht

Maak een screenshot waarop je de uiteindelijke situatie laat zien.

Je screenshot moet minimaal bevatten:

```bash
ls -ld /srv/teams/support
ls -ld /srv/teams/development
```

Daarnaast moet je aantonen dat:

1. twee gebruikers van hetzelfde team een bestand kunnen gebruiken;
2. een gebruiker van een ander team geen toegang heeft;
3. het setgid-bit actief is.

---

# Inleveren

Lever één document in met:

* je naam;
* de gebruikte Linux-distributie;
* de gebruikte commando's;
* screenshots van de belangrijkste stappen;
* antwoorden op vraag 1 t/m 10;
* screenshots van de eindtest.

**Bestandsnaam:**

```text
Linux_Gebruikers_Groepen_Rechten_Voornaam_Achternaam.pdf
```

**Tijdsduur:** maximaal 1,5 uur.

